%% EAAI submission draft -- elsarticle (author-year / elsarticle-harv).
%% Assembled 2026-08-18 from paper/sections/*.md (all FROZEN versions).
%% Compile: pdflatex -> bibtex -> pdflatex x2  (or upload this folder to Overleaf).
%% TODO before submission: author affiliation; move Highlights to the
%% separate submission field (kept in highlights.tex).

\documentclass[preprint,12pt,authoryear]{elsarticle}

\usepackage{graphicx}
\usepackage{booktabs}
\usepackage{amsmath}
\usepackage{url}

\newcommand{\mapm}{mAP\(_{50\text{--}95}\)}

\journal{Engineering Applications of Artificial Intelligence}

\begin{document}

\begin{frontmatter}

\title{Allocation-aware inpainting augmentation for object detection:
the class set, not the per-class weighting, determines where
performance gains occur}

\author[a]{Daehyun Yoo}
\ead{dhyoo970111@gmail.com}
\affiliation[a]{organization={Affiliation to be completed},
                city={},
                country={Republic of Korea}}

\begin{abstract}
Synthetic images produced by diffusion inpainting can enlarge
object-detection training sets at low cost, but a practical question
remains open: given a fixed budget of generated images, which classes
should receive them? The contribution in artificial intelligence is a
controlled answer to this allocation question. We cross the allocation
signal --- a frequency-defined tail set versus a difficulty-defined weak
set selected by measured per-class average precision --- with the
within-set weighting (uniform versus priority-weighted), holding the
budget, class count, generator, prompts, and rule-based quality
verification fixed, and we evaluate the four resulting augmentation arms
with pre-registered, hash-frozen contrasts, seed-blocked tests, and
multiple-comparison correction. The engineering application is military
aircraft detection in two operationally distinct regimes: a 43-class
ground-level natural-view benchmark and a 20-class satellite-view
benchmark, each paired with a compact one-stage detector and a large
pretrained transformer detector. Three results follow. First, the
allocation signal determines where gains occur: each arm improves the
classes it targets (set-by-scope interaction $+0.06$, corrected
$p = 0.017$; replicated at $+0.05$ on the satellite benchmark), while
the weighting axis is null in all four comparisons. Second, every
augmentation gain vanishes under the stronger detector, whose
unaugmented baseline exceeds the best augmented compact model. Third,
zero-cost repeat-factor resampling dominates generation on the deeply
imbalanced benchmark but is null on the shallow one, where inpainting is
the only augmentation with significant gains. These findings yield a
three-step engineering rule --- check detector headroom, probe with free
resampling, then generate for the measured-weakest classes with a
uniform split. Code, generated image pools, and all training artifacts
will be released.
\end{abstract}

\begin{keyword}
Data augmentation \sep Diffusion inpainting \sep Object detection \sep
Class imbalance \sep Synthetic data allocation \sep Military aircraft
detection
\end{keyword}

\end{frontmatter}

\section{Introduction}\label{sec:intro}

Object detection systems deployed in specialized domains face a chronic
data problem: collecting and annotating images for every class of
interest is expensive, and the resulting training sets are almost
inevitably class-imbalanced. Military aircraft recognition is a
representative case --- new airframes appear rarely in public imagery,
per-class collection costs vary by orders of magnitude, and yet the
operational cost of missing a rare type is high. Diffusion-based image
generation has made synthetic training data cheap enough to be a
practical remedy, and a growing literature shows that generated images
can improve detectors. For the practicing engineer, however, the binding
question is no longer \emph{whether} to generate but \emph{where to
spend a fixed generation budget}: given the capacity to produce, say,
one thousand verified synthetic images, which classes should receive
them?

The long-tail literature supplies a default answer: give them to the
rare classes. This prescription rests on a premise so common that it is
rarely stated --- that training frequency identifies the classes that
need help. On our natural-view benchmark the premise fails. Per-class
average precision (AP) of a strong baseline correlates \emph{negatively}
with training frequency (Pearson $r = -0.375$): the rarest aircraft
types are often visually distinctive and easy, while the hardest classes
--- visually confusable fighter jets --- sit in the middle and head of
the frequency distribution. Frequency and measured difficulty therefore
select \emph{disjoint} sets of classes to support, and ``help the
rare'' and ``help the weak'' become competing allocation policies rather
than synonyms. Which signal should govern the budget, and does the
choice matter at all?

We answer this with a controlled factorial experiment rather than
another method. Using background inpainting as the generation mechanism
--- annotated objects are geometrically protected and only their
surroundings are regenerated, so labels remain exact by construction ---
we cross the \emph{allocation signal} (frequency-defined tail set vs.\
AP-defined weak set) with the \emph{within-set weighting} (uniform vs.\
priority-weighted) while holding the budget, the number of target
classes, the generator, the prompts, and the quality-verification
protocol fixed. The four resulting augmentation arms differ in exactly
one respect: how the same budget is allocated. We run this design on two
public benchmarks in operationally distinct imaging regimes --- a
43-class natural-view dataset and the 20-class satellite-view MAR20 ---
and under two detectors that differ by an order of magnitude in
capacity, evaluating all arms with pre-registered contrasts, seed-level
statistics, and multiple-comparison correction.

Three findings emerge. First, \emph{the allocation signal determines
where performance gains occur}: each arm improves the class set it
targets by several points of AP while conferring only a smaller
spillover elsewhere, a double dissociation confirmed by a pre-registered
interaction test on the natural-view dataset and replicated in effect
size on MAR20. In contrast, the within-set weighting axis is null in all
four comparisons --- \emph{which classes} receive the budget matters;
\emph{how many images each class gets} does not measurably matter at the
budgets we study. Second, every augmentation effect we measure is
conditional on detector headroom: under a large pretrained transformer
detector whose baseline already exceeds the best augmented compact
model, all gains vanish. Third, the zero-cost resampling baseline ---
repeat-factor sampling --- dominates every generative arm on the deeply
imbalanced dataset yet is null on the shallowly imbalanced one, where
inpainting is the only augmentation with significant gains; resampling's
outcome thus serves as a cheap diagnostic of whether generation is worth
its cost. Together these results yield a three-step decision rule for
practitioners --- check headroom, probe with free resampling, then
generate for the \emph{measured-weakest} classes with a uniform split
--- and a caution for the literature: augmentation studies that vary a
single recipe on a single detector risk reporting effects that a
stronger baseline model would erase.

Our contributions are:
\begin{itemize}
\item \textbf{An allocation-signal finding.} Under a fixed budget, fixed
  class count, and fixed generation pipeline, the choice of class set
  --- not the per-class weighting --- determines where
  inpainting-augmentation gains occur (set-by-scope interaction $+0.05$
  to $+0.06$ across two datasets; weighting null in all four
  pre-registered comparisons), and frequency is a poor default signal
  wherever frequency and difficulty decouple.
\item \textbf{An evaluation protocol for augmentation claims.} A
  $2\times2$ factorial with budget control, pre-registered and
  hash-frozen contrasts, seed-blocked tests with Holm correction, and
  rule-based generation verification --- a reusable template that blocks
  the confounds (budget differences, single-seed effects, post-hoc
  hypothesis selection) common in augmentation studies.
\item \textbf{A condition map for generative augmentation.} Evidence
  that augmentation gains are moderated by detector headroom and by the
  depth of dataset imbalance, including the failure modes: all gains
  vanish under a strong pretrained detector, and resampling turns
  harmful on the weak scope of a shallowly imbalanced dataset.
\item \textbf{An open, deterministic pipeline.} Bbox-protected
  inpainting with quantitative quality verification, hallucination
  auditing, and fully seed-deterministic regeneration, released with
  plans, logs, metrics, and weights for all 66 training runs.
\end{itemize}

The remainder of the paper reviews related work
(Section~\ref{sec:related}), details the method and experimental setup
(Sections~\ref{sec:method}--\ref{sec:setup}), reports results
(Section~\ref{sec:results}), develops engineering guidance, a
mechanistic interpretation, and limitations
(Section~\ref{sec:discussion}), and concludes
(Section~\ref{sec:conclusion}).

\section{Related work}\label{sec:related}

\subsection{Generative augmentation for object detection}

Synthetic training data for detection has largely followed a
\emph{foreground} line of work: Copy-Paste showed that pasting real
object instances into new images is a strong augmentation on its own
\citep{ghiasi2021copy}; X-Paste scaled the recipe by generating the
pasted instances with text-to-image diffusion \citep{zhao2023xpaste};
MosaicFusion synthesizes rare-category objects together with their masks
\citep{xie2024mosaicfusion}; Gen2Det generates scene-centric images by
grounded inpainting and couples them with image- and instance-level
filtering \citep{suri2023gen2det}; and DiverGen attributes much of the
benefit to the \emph{diversity} of generated data \citep{fan2024divergen}.
Domain-specific variants fine-tune or condition the generator for a
target domain, as in ODGEN for specialized detection datasets
\citep{zhu2024odgen} and AeroGen for layout-conditioned remote-sensing
imagery \citep{tang2025aerogen}. Closest to our pipeline is the
\emph{background} line: \citet{li2024background} protect annotated
foregrounds and regenerate only the surroundings with Stable Diffusion
inpainting, reporting that background augmentation can outperform
object-centric augmentation for robustness and generalization. We adopt
this mechanism deliberately --- it keeps labels exact by construction
--- but ask a question this literature leaves open. Existing studies
establish \emph{that} generated data can help and how to filter it; the
budget itself is typically spent uniformly or left as an unexamined
implementation choice. None isolates the \emph{allocation signal} ---
which classes receive the budget --- as the experimental variable, which
is precisely the factor our design manipulates under a fixed budget,
generator, and quality-verification protocol.

\subsection{Long-tailed detection: resampling and rebalancing}

The long-tail literature defines the problem by class frequency. LVIS
introduced the large-vocabulary benchmark, the rare/common/frequent
reporting convention, and repeat-factor sampling (RFS), the standard
image-level resampling scheme \citep{gupta2019lvis}; instance-aware and
exponentially weighted refinements followed \citep{yaman2023irfs,
ahmed2025eirfs}. Loss-side rebalancing spans balanced group softmax
\citep{li2020bags}, Seesaw loss \citep{wang2021seesaw}, and equalized
focal loss for one-stage detectors \citep{li2022efl}. Two findings from
this literature shape our protocol. First, rebalancing techniques
developed on two-stage detectors transfer unevenly to modern one-stage
pipelines, where strong default augmentation (mosaic, mixup) already
dominates sampling- and loss-based corrections \citep{crasto2024imbalance};
accordingly, all of our arms are evaluated as \emph{marginal} additions
on top of a strong augmented baseline rather than in isolation --- the
reporting practice of X-Paste and Gen2Det. Second, and central to our
motivation, the entire resampling tradition presumes that frequency
identifies the classes that need help. Our data violate that premise ---
per-class average precision correlates \emph{negatively} with training
frequency on the natural-view dataset --- and our factorial design
treats the frequency-defined tail and the measured-AP-defined weak set
as competing allocation signals rather than assuming either one.

\subsection{Allocating and selecting synthetic data}

A smaller body of work asks not whether to generate but \emph{where} and
\emph{what}. Diffusion-curriculum approaches allocate synthetic data
adaptively along an image-guidance spectrum for tail classes
\citep{liang2025discl}; and generating informative samples with
class-aware generative adversarial networks for weak regions of a
classifier was explored as early as the GAN era, in medical image
classification \citep{bozorgtabar2019informative}, anticipating our
difficulty signal. On the selection side, quality filtering by CLIP
score or detector feedback is by now a standard component
\citep{suri2023gen2det, fan2024divergen}, and recent work cautions that
retrieved real images can match or outperform synthetic ones
\citep{geng2024unmet}, reinforcing the need for strong non-generative
baselines. These works motivate difficulty-aware allocation but evaluate
it entangled with other design changes --- curricula, generator
fine-tuning, filtering --- and without budget-controlled comparisons.
Our contribution to this strand is methodological: a $2\times2$
factorial that crosses the allocation signal with the within-set
weighting under an identical budget, class count, generator, and
verification protocol, evaluated against pre-registered contrasts, so
that the effect of the signal itself is identified.

\subsection{Aircraft detection in surveillance imagery}

Military aircraft recognition is served by dedicated benchmarks ---
notably MAR20, the largest public military-aircraft remote-sensing
dataset \citep{yu2023mar20} --- and by a line of detector-architecture
work, from fine-grained one-stage designs evaluated on MAR20
\citep[FGA-YOLO;][]{wu2025fga} to small-object detection for aerial
platforms published in this journal \citep{song2024small}. This
literature concentrates on architectural improvements under fixed data;
the complementary engineering question --- how to spend a limited
synthetic-data budget when per-class collection is expensive --- has
received little controlled study in the domain. We address that question
on both a ground-level natural-view benchmark and MAR20, and use the two
detector families of Section~\ref{sec:detectors} to test whether the
answer survives a change of architecture and capacity.

\section{Method}\label{sec:method}

\subsection{Overview}\label{sec:overview}

Our pipeline augments a detection training set with synthetic images in
which the \emph{background} is regenerated by a diffusion inpainting
model while every annotated object is geometrically protected
(Figure~\ref{fig:pipeline}). Given a fixed budget of $B$ synthetic
images and a fixed number of target classes $K$, an \emph{allocation
plan} assigns a per-class image quota; training images containing the
targeted class are drawn as sources, inpainted outside all annotated
object boxes, required to pass rule-based quality verification against
quantitative criteria (Section~\ref{sec:verify}), and inserted into the
training split only. The design isolates a single question: holding the
generator, prompts, quality verification, budget, and class count fixed,
does the \emph{allocation signal} --- the rule that selects which
classes receive the budget --- change where performance gains occur?

\begin{figure}[t]
\centering
\includegraphics[width=\linewidth]{figures/fig1_pipeline_design.pdf}
\caption{Overview. (a)~Generation pipeline: all annotated boxes in a
source image are geometrically protected; diffusion inpainting
regenerates only the background; outputs must pass rule-based
verification (with budget refill) before insertion into the training
split. Budget $B$, class count $K$, generator, prompts, and verification
are held fixed across arms. (b)~The $2\times2$ allocation design: class
set (frequency-defined tail vs.\ AP-defined weak) $\times$ within-set
weighting (uniform vs.\ priority-weighted).}
\label{fig:pipeline}
\end{figure}

\subsection{Allocation design: class set $\times$ within-set weighting}\label{sec:alloc}

We instantiate two allocation signals and two within-set weightings,
yielding a $2\times2$ factorial over augmentation arms:
\begin{itemize}
\item \textbf{Tail set (frequency signal).} The bottom 30\% of classes
  by training instance count --- the conventional long-tail
  prescription.
\item \textbf{Weak set (difficulty signal).} The $K$ classes with the
  lowest baseline per-class average precision (AP), measured on the
  validation split by the \texttt{basic\_aug} baseline detector and
  averaged over its training seeds. $K$ is fixed to the tail-set size so
  that set identity, not set size, is the manipulated variable.
\item \textbf{Uniform weighting} splits $B$ equally across the $K$
  classes.
\item \textbf{Weighted allocation} distributes $B$ proportionally to a
  priority score: for the tail set, a convex combination of rarity and
  measured weakness (weight 0.6 toward rarity); for the weak set,
  measured weakness alone.
\end{itemize}

Quotas are computed by a deterministic capped largest-remainder
allocator: each class first receives a minimum quota, the remainder is
apportioned proportionally to the weighting under per-class caps, and
fractional leftovers go to the largest remainders with index-ordered tie
breaking. The allocator guarantees that every arm spends exactly $B$
images on exactly $K$ classes, so budget can never confound the
comparison. Empirically, the two sets are disjoint on the natural-view
benchmark and overlap only marginally on the satellite-view benchmark
(Section~\ref{sec:datasets}) --- frequency and measured difficulty
select genuinely different classes here (Section~\ref{sec:r1}), which is
what renders the factorial informative.

\subsection{Background inpainting generation}\label{sec:generation}

Synthetic images are produced with Stable Diffusion inpainting
\citep[\texttt{runwayml/stable-diffusion-inpainting};][]{rombach2022ldm}
at 16-bit floating-point precision, 20 denoising steps, guidance scale
7.5, and strength 0.85, at $512\times512$ working resolution and resized
back to the source resolution. For each source image we build a mask
that \emph{protects} all annotated boxes with a 10\% padding margin and
a blurred boundary band, so that diffusion may alter only the
background. Scene prompts are dataset-specific (five natural-view
aviation prompts for MAD; five top-down aerodrome prompts for MAR20), as
are the negative prompts, which suppress --- among other artifacts ---
the generation of new aircraft. After generation, the protected regions
are pasted back from the source pixels, making object appearance and
labels exact by construction; source labels are copied unchanged.
Generation is fully deterministic given the plan: each output is a
function of the source image, the prompt rotation, and a derived
per-image seed, which enables byte-identical regeneration and safe reuse
across experiments (examples in Figure~\ref{fig:examples}).

\subsection{Rule-based quality verification and budget refill}\label{sec:verify}

Source images must first offer enough editable background for inpainting
to be meaningful (protected regions covering no more than 95\% of the
frame); sources failing this pre-check are excluded before generation.
Every generated image must then pass two quantitative checks: (i)~the
background changed substantially (mean absolute pixel difference $\geq
10$ gray levels outside the boxes), and (ii)~the protected interiors did
not change beyond a compression-noise tolerance ($\leq 5$ gray levels).
Failed generations are retried and, if still failing, replaced from
alternative sources within a bounded attempt budget until the class
quota is met, so all arms reach the same \emph{realized} budget. A run
aborts if more than half of all attempts fail, flagging generator
malfunction. As an additional audit --- reported as a robustness check,
not a filter --- a detector-based scan counts \emph{hallucinated}
objects that diffusion occasionally paints into backgrounds, and an
object-level gate variant retrains the arms with all affected images
removed (Section~\ref{sec:robust}).

\subsection{Training arms and baselines}\label{sec:arms}

All arms share the detector, schedule, and the default
photometric/geometric augmentation of the Ultralytics training framework
\citep{jocher2023ultralytics}; they differ only in training data.
\texttt{real\_only} disables framework augmentation (reference lower
bound); \texttt{basic\_aug} is the primary baseline that all gains are
measured against. Repeat-factor sampling (RFS) is the zero-cost
resampling baseline: with $f(c)$ the fraction of training images
containing class $c$, each image $I$ is duplicated according to
$r(I) = \max_{c \in I} r(c)$, where
$r(c) = \max\!\bigl(1, \sqrt{t/f(c)}\bigr)$ with $t = 0.1$, materialized
at the dataset level. The four inpainting arms add their $B$-image pool
to the training split; validation and test splits are never touched.

\subsection{Statistical protocol}\label{sec:stats}

We separate planning from evaluation. All plan-forming measurements
(weak set selection, priority scores) use the validation split only; the
test split is reserved for final reporting, and no plan or
hyperparameter was revised after observing test metrics. For the
full-factorial cells, five contrasts were specified and frozen --- with
a content hash and timestamp --- before any test-split metric existed:
(1)~tail arms vs.\ baseline on the tail scope, (2)~weak arms vs.\
baseline on the weak scope, (3)~the set-by-scope interaction,
(4--5)~weighted vs.\ uniform within each set. Contrasts are evaluated on
per-seed macro AP (the unweighted mean of per-class AP over the classes
in a scope) with seed-blocked paired $t$-tests ($n = 3$ independent
training seeds) and Holm correction \citep{holm1979simple} within each
cell's five-contrast family; all tests are two-sided at $\alpha = 0.05$.
Wilcoxon signed-rank tests on seed-averaged per-class AP provide
secondary, distribution-free evidence. We report evaluation over three
class scopes: \emph{all}, the frequency-defined \emph{tail}, and the
AP-defined \emph{weak} set.

\section{Experimental setup}\label{sec:setup}

The engineering application is military aircraft detection in two
operationally distinct imaging regimes --- ground-level natural views
and satellite views --- a setting where per-class data collection is
expensive and class imbalance is unavoidable, making augmentation budget
decisions a practical engineering concern. Both benchmarks are publicly
available, supporting replication of every result.

\subsection{Datasets}\label{sec:datasets}

\textbf{MAD} \citep[Military Aircraft Detection, natural side/ground
views; public;][]{rookieengg2023mad} contains 11{,}788 images
(9{,}430/1{,}179/1{,}179 train/val/test) with 18{,}832 boxes over 43
aircraft types; the train-split imbalance ratio is 10.5. \textbf{MAR20}
\citep[Military Aircraft Recognition dataset, satellite view, horizontal
boxes; public;][]{yu2023mar20} provides an official test split of
2{,}511 images, which we preserve untouched for final reporting; the
official 1{,}331 training images are split 80/20 into train/val with
class-stratified sampling. MAR20's 20-class distribution is markedly
shallower than MAD's, which Section~\ref{sec:r3} exploits as a natural
contrast in imbalance structure. The tail set is a property of the
dataset; the weak set is derived from each detector's own baseline
(Section~\ref{sec:alloc}) and is therefore cell-specific. On MAD the two
sets contain 13 classes each and are disjoint for both detectors; on
MAR20 they contain 6 classes each, overlapping in one class for the
compact detector and two for the transformer.

\subsection{Detectors and training}\label{sec:detectors}

We use two detectors spanning an order of magnitude in capacity:
YOLOv8n, a compact one-stage detector of the You-Only-Look-Once (YOLO)
family (3.2M parameters), and RT-DETR-L, the large variant of the
Real-Time Detection Transformer \citep[32M parameters;][]{zhao2024rtdetr}.
Both are initialized from checkpoints pretrained on the Common Objects
in Context (COCO) benchmark, so the cells differ in capacity and
architecture family, not in pretraining corpus. Both are trained at
$640\times640$ for 50 epochs with early stopping (patience 15) under
identical schedules across arms; the batch size is fixed to 8 for
RT-DETR-L and set automatically for YOLOv8n, with all remaining
hyperparameters at framework defaults. Every arm--cell combination is
trained with three independent seeds; all reported statistics aggregate
over seeds as described in Section~\ref{sec:stats}.

\subsection{Cells and budgets}\label{sec:cells}

Table~\ref{tab:cells} summarizes the four dataset $\times$ detector
cells. The YOLOv8n cells run the full $2\times2$ factorial plus
baselines; the RT-DETR-L cells run a reduced design (baseline, RFS, and
the two tail-set arms) intended to test whether gains transfer to a
stronger detector rather than to re-estimate the interaction. The
uniform arm's pool depends only on the dataset and is shared across the
detector cells; the weighted arm's difficulty component is recomputed
from each detector's own baseline validation AP, so its allocation is
detector-specific. Budgets scale with dataset size: $B = 1{,}000$
($K = 13$, per-class quota 5--200) for MAD and $B = 500$ ($K = 6$, quota
5--100) for MAR20. Realized budgets equal planned budgets in all cells
(verification failure rates: $\leq 21\%$ of attempts on MAD, 0.6\% on
MAR20, absorbed by the refill mechanism of Section~\ref{sec:verify}).

\begin{table}[t]
\centering
\caption{Experimental cells. All cells share the generation,
verification, and statistical protocol.}
\label{tab:cells}
\begin{tabular}{llll}
\toprule
Dataset & Detector & Arms & Budget $B$\,/\,$K$ \\
\midrule
MAD   & YOLOv8n   & full $2\times2$ + RFS + baselines & 1{,}000\,/\,13 \\
MAD   & RT-DETR-L & reduced (tail-set arms + RFS)     & 1{,}000\,/\,13 \\
MAR20 & YOLOv8n   & full $2\times2$ + RFS + baselines & 500\,/\,6 \\
MAR20 & RT-DETR-L & reduced (tail-set arms + RFS)     & 500\,/\,6 \\
\bottomrule
\end{tabular}
\end{table}

\subsection{Reproducibility}

Generation is seed-deterministic given a plan
(Section~\ref{sec:generation}); plans, generation logs, verification
reports, contrast freezes (with content hashes), per-class metrics, and
trained weights for all 66 primary training runs --- plus the
object-gate retrainings of Section~\ref{sec:robust} --- are archived,
and the full pipeline, from dataset normalization to the statistical
reports, runs from configuration files without manual intervention.
Code and artifacts will be released upon publication.

\section{Results}\label{sec:results}

We report all results on the military aircraft detection cells of
Section~\ref{sec:setup} as test-split mean average precision over
intersection-over-union thresholds 0.50--0.95 (\mapm), and give the mean
over three independent training seeds with per-arm standard deviations
computed across seeds. Gains ($\Delta$) are always relative to the
\texttt{basic\_aug} baseline of the same cell. Pre-registered contrasts
(Section~\ref{sec:stats}) are evaluated with seed-blocked paired
$t$-tests, 95\% confidence intervals (CIs), and Holm correction within
each cell's five-contrast family; class-level Wilcoxon signed-rank tests
on seed-averaged per-class AP serve as secondary evidence.
Table~\ref{tab:master} gives the master results for all four cells;
Table~\ref{tab:contrasts} lists the pre-registered contrast outcomes.

\subsection{The allocation signal determines where gains occur}\label{sec:r1}

\paragraph{Frequency and difficulty select different classes}
On MAD, per-class training-instance count and baseline test AP correlate
\emph{negatively} (\mapm{} basis; Pearson $r = -0.375$, $p = 0.013$;
Spearman $\rho = -0.408$, $p = 0.007$; Figure~\ref{fig:freqap}). In
AP\(_{50}\) terms, the rarest class (YF-23, 97 training instances)
attains the highest score (0.93), whereas the weakest classes (Rafale
0.31, F-14 0.37) sit in the medium and head frequency ranges.
Empirically, the frequency-defined tail set and the AP-defined weak set
are disjoint on MAD (0/13 overlap) and overlap in a single class on
MAR20 for the compact detector (Section~\ref{sec:datasets}). This is the
premise that makes the $2\times2$ allocation design informative: if
frequency were a good proxy for difficulty, the two signals would select
the same classes and could not be dissociated.

\begin{figure}[t]
\centering
\includegraphics[width=0.85\linewidth]{figures/fig2_freq_vs_ap.pdf}
\caption{Frequency is a poor proxy for difficulty on MAD. Per-class
baseline test \mapm{} (three-seed mean) against training instance count
(log scale). The frequency-defined tail set (blue) and the AP-defined
weak set (orange) are disjoint; the rarest class (YF-23) is the easiest,
while the weakest classes (Rafale, F-14) sit at medium-to-head
frequency.}
\label{fig:freqap}
\end{figure}

\paragraph{Gains land on the targeted set (MAD, confirmatory)}
Figure~\ref{fig:dissoc} (left) shows the result on MAD with YOLOv8n as a
double dissociation: each set family improves its own target scope
substantially more than the other family does, and vice versa. Averaged
over seeds, the two tail-set arms improved the tail scope by $+0.061$
over baseline while improving the weak scope by only $+0.019$; the two
weak-set arms showed the mirrored pattern ($+0.040$ on weak, $+0.022$ on
tail). The pre-registered set-by-scope interaction --- the
difference-in-differences of the two set families across the two scopes
--- was $+0.0598$ (95\% CI $[0.045, 0.075]$, Holm-adjusted
$p = 0.017$), with all three seeds showing the same direction. At the
class level, gains were broad rather than concentrated: 11 of 13 tail
classes improved under the tail-uniform arm (largest: Tu-95 $+0.169$,
MiG-31 $+0.101$), and 11 of 13 weak classes improved under the
weak-uniform arm (largest: Tornado $+0.123$, Rafale $+0.086$ --- the
weakest class in the dataset).

\begin{figure}[t]
\centering
\includegraphics[width=\linewidth]{figures/fig3_dissociation.pdf}
\caption{The allocation signal determines where gains occur. $\Delta$
\mapm{} versus the same-cell baseline for the four inpainting arms on
the tail scope (blue) and weak scope (orange); black-edged bars mark each arm's targeted scope; error bars span
$\pm 1$ standard deviation over three seeds. Left: MAD
(confirmatory; interaction Holm-adjusted $p = 0.017$). Right: MAR20
(replication of the effect size; Section~\ref{sec:r1}).}
\label{fig:dissoc}
\end{figure}

\paragraph{The dissociation replicates in the satellite domain (MAR20)}
Figure~\ref{fig:dissoc} (right) shows the same analysis for MAR20 with
YOLOv8n. The interaction estimate was $+0.0523$ (95\% CI
$[0.021, 0.083]$; positive in all three seeds), which lies inside the
95\% CI of the MAD estimate; the tail-targeted arm again gained most on
the tail scope (selective: $+0.044$; class-level Wilcoxon $p = 0.031$,
the smallest attainable value at $n = 6$) and the weak-targeted arms
most on the weak scope ($+0.030$ to $+0.034$). Because the MAR20 scopes
contain only six classes each, the seed-blocked tests are underpowered
and the interaction does not survive Holm correction (unadjusted
$p = 0.018$, Holm-adjusted $p = 0.091$); we therefore report MAR20 as an
independent replication of the effect size rather than a second
confirmatory result.

\paragraph{Within-set weighting shows no benefit}
In contrast to the choice of class set, \emph{how} the fixed budget is
divided inside the chosen set showed no evidence of benefit in any
condition: the weighted-versus-uniform contrasts were null in all four
pre-registered comparisons, with estimates centred near zero ($+0.002$,
$p = 0.81$ and $-0.012$, $p = 0.17$ on MAD; $+0.014$, $p = 0.11$ and
$+0.004$, $p = 0.52$ on MAR20; Table~\ref{tab:contrasts}). The CIs do
not exclude small effects in either direction, so we do not claim strict
equivalence; what the data rule out is a within-set weighting benefit
comparable to the set-choice effect. Given that the uniform and weighted
arms shared the same class sets, generator, prompts, quality
verification, and budget, we conclude that at the budgets studied the
allocation \emph{signal} (which classes) --- not the allocation
\emph{shape} (how many per class) --- is the operative variable.

\paragraph{Targeted gains come with a smaller global spillover}
Augmenting one set also produced consistent but smaller gains outside
it: on MAD the tail-uniform arm improved the weak scope by $+0.031$
(class-level $p = 0.002$), and on MAR20 the cross-set gains ranged from
$+0.004$ to $+0.013$. Background inpainting therefore appears to combine
a set-specific targeting component with a smaller domain-wide component,
consistent with the recall-oriented mechanism discussed in
Section~\ref{sec:discussion}.

\subsection{A strong pretrained detector removes all augmentation gains}\label{sec:r2}

Replacing the compact detector with RT-DETR-L
(Section~\ref{sec:detectors}) raised the baseline from 0.575 to 0.869 on
MAD and from 0.572 to 0.697 on MAR20 --- and eliminated every
augmentation effect we measured (Figure~\ref{fig:condmap},
Table~\ref{tab:master}). On MAD, all deltas were zero to negative:
uniform $-0.004$, selective $-0.001$, and RFS $-0.016$; the pooled
inpainting-versus-baseline contrast on the tail scope was $-0.008$
($p = 0.36$). On MAR20 the pattern repeated (uniform $-0.007$, selective
$-0.005$, RFS $+0.003$; pooled tail contrast $+0.004$, $p = 0.72$).
Notably, the strong detector's \emph{baseline} already exceeded the best
augmented YOLOv8n result by a wide margin (0.869 vs.\ 0.661 on MAD).

Two observations follow. First, augmentation gains in this study are
conditional on detector headroom: once a model retains little headroom
on the dataset, neither resampling nor generative augmentation adds
measurable value, and the largest single lever is the detector upgrade
itself (a cross-cell comparison; Section~\ref{sec:rules}). Second, the
only consistently \emph{negative} delta under the strong detector was
RFS on MAD ($-0.016$ on all three scopes), suggesting that increasing
duplicate exposure can be mildly harmful when capacity and pretraining
already suffice; we report this as an observation without claiming a
mechanism. Because the RT-DETR-L cells implemented a reduced design
(baseline, RFS, and the two tail-set arms only), no interaction test is
available there; the ``vanishing gains'' claim rests on the near-zero
deltas of all measured arms.

\begin{figure}[t]
\centering
\includegraphics[width=0.9\linewidth]{figures/fig4_condition_map.pdf}
\caption{Condition map across the four cells. Best augmentation
$\Delta$ \mapm{} (all scope) against the unaugmented baseline of each
cell; triangles denote repeat-factor sampling, circles the best
inpainting arm; color denotes the dataset. Gains shrink to zero under
the high-baseline RT-DETR-L cells (right), and resampling and inpainting
exchange ranks between MAD and MAR20 (left).}
\label{fig:condmap}
\end{figure}

\subsection{Resampling effectiveness depends on the imbalance structure}\label{sec:r3}

Repeat-factor sampling --- the zero-cost resampling baseline that
dominated every diffusion arm on MAD ($+0.086$ all, $+0.098$ tail,
$+0.117$ weak; class-level $p \leq 2\times10^{-4}$) --- transferred
poorly to MAR20. With the same detector and protocol, RFS on MAR20
yielded $+0.003$ (all), $+0.016$ (tail), and $-0.011$ (weak;
class-level $p = 0.031$) --- the only statistically significant negative
augmentation effect we observed under YOLOv8n. MAD is deeply imbalanced
across 43 classes (max/min instance ratio 10.5), so re-exposing
rare-class images carries additional training signal; MAR20's 20-class
distribution is comparatively shallow, and additional exposure of
already-seen images adds little, while inpainting --- which injects
genuinely new background variation --- remains effective
(Section~\ref{sec:r1}). On MAR20, inpainting was in fact the only
augmentation with statistically significant gains. The practical
implication is developed in Section~\ref{sec:discussion}: RFS is a cheap
first attempt rather than a universal winner, and its failure is
diagnostic of a regime where generative augmentation is worth its cost.

\subsection{Robustness checks}\label{sec:robust}

\paragraph{Generation quality does not explain the dissociation}
The three originally generated MAD pools passed rule-based quality
verification at similar rates (79--82\%), and the fourth pool reached
its quota under the same criteria; per-pool Fr\'{e}chet inception
distance \citep[FID;][]{heusel2017fid} and CLIP-based image--text
similarity scores \citep{radford2021clip} were comparable, with the
\emph{worst}-FID pool (weakness) producing the significant weak-scope
gain --- quality differences cannot account for where gains landed. On
MAR20 the verification failure rate was 0.6\%, and all four pools
reached the full 500-image budget.

\paragraph{Hallucinated objects do not drive the RFS gap on MAD}
A detector-based audit found unlabeled aircraft hallucinated into 10.2\%
of MAD synthetic backgrounds (see Figure~\ref{fig:examples}, top row).
Removing all affected images with an object-level gate and retraining
the three original diffusion arms changed their all-scope performance by
at most $|\Delta| \leq 0.017$ (all nine gated-versus-ungated comparisons
non-significant, $p = 0.11$--$0.97$), while the RFS advantage persisted
essentially unchanged ($\Delta = -0.055$ to $-0.067$ versus RFS,
$p < 10^{-4}$). Label noise from generation artifacts therefore explains
neither the dissociation nor the resampling gap.

\begin{figure}[t]
\centering
\includegraphics[width=0.8\linewidth]{figures/fig5_generation_examples.pdf}
\caption{Generation examples (source, protection mask, output; mask
black = protected). Rows: two MAD and two MAR20 cases spanning
formation, airshow, apron, and airbase scenes. The top row also
illustrates the hallucination failure mode quantified in
Section~\ref{sec:robust} --- diffusion has painted additional, unlabeled
aircraft into the regenerated sky --- which motivates the detector-based
audit and the object-gate robustness check.}
\label{fig:examples}
\end{figure}

\begin{table}[t]
\centering
\caption{Master results: $\Delta$ \mapm{} versus the same-cell
\texttt{basic\_aug} baseline (mean over three seeds; baseline rows show
absolute mAP $\pm$ SD). Bold marks the scope targeted by each arm;
scopes are cell-specific (Section~\ref{sec:datasets}).}
\label{tab:master}
\small
\begin{tabular}{llrrr}
\toprule
Dataset $\times$ detector & Arm & all & tail & weak \\
\midrule
MAD $\times$ YOLOv8n
  & basic\_aug    & $0.575\pm0.007$ & $0.620\pm0.009$ & $0.456\pm0.030$ \\
  & RFS           & $+0.086$ & $+0.098$ & $+0.117$ \\
  & tail-uniform  & $+0.031$ & $\mathbf{+0.060}$ & $+0.031$ \\
  & tail-weighted & $+0.023$ & $\mathbf{+0.062}$ & $+0.007$ \\
  & weak-uniform  & $+0.016$ & $+0.014$ & $\mathbf{+0.046}$ \\
  & weak-weighted & $+0.020$ & $+0.030$ & $\mathbf{+0.034}$ \\
\midrule
MAD $\times$ RT-DETR-L
  & basic\_aug    & $0.869\pm0.006$ & $0.889\pm0.008$ & $0.842\pm0.012$ \\
  & RFS           & $-0.016$ & $-0.010$ & $-0.016$ \\
  & tail-uniform  & $-0.004$ & $\mathbf{-0.013}$ & $-0.006$ \\
  & tail-weighted & $-0.001$ & $\mathbf{-0.004}$ & $+0.003$ \\
\midrule
MAR20 $\times$ YOLOv8n
  & basic\_aug    & $0.572\pm0.005$ & $0.532\pm0.013$ & $0.467\pm0.017$ \\
  & RFS           & $+0.003$ & $+0.016$ & $-0.011$ \\
  & tail-uniform  & $+0.010$ & $\mathbf{+0.030}$ & $+0.013$ \\
  & tail-weighted & $+0.012$ & $\mathbf{+0.044}$ & $+0.008$ \\
  & weak-uniform  & $+0.004$ & $+0.008$ & $\mathbf{+0.030}$ \\
  & weak-weighted & $+0.003$ & $+0.004$ & $\mathbf{+0.034}$ \\
\midrule
MAR20 $\times$ RT-DETR-L
  & basic\_aug    & $0.697\pm0.001$ & $0.676\pm0.007$ & $0.611\pm0.010$ \\
  & RFS           & $+0.003$ & $+0.013$ & $0.000$ \\
  & tail-uniform  & $-0.007$ & $\mathbf{+0.007}$ & $-0.009$ \\
  & tail-weighted & $-0.005$ & $\mathbf{0.000}$  & $-0.007$ \\
\bottomrule
\end{tabular}
\end{table}

\begin{table}[t]
\centering
\caption{Pre-registered contrasts (seed-blocked paired $t$, $n = 3$
seeds; Holm correction within each cell's five-contrast family). Both
full-factorial cells use YOLOv8n (Section~\ref{sec:cells}).}
\label{tab:contrasts}
\small
\begin{tabular}{llrrrr}
\toprule
Dataset & Contrast & Estimate & 95\% CI & $p$ & Holm $p$ \\
\midrule
MAD & Tail arms vs.\ baseline (tail scope)  & $+0.061$ & $[+0.026, +0.096]$ & 0.017 & 0.069 \\
MAD & Weak arms vs.\ baseline (weak scope)  & $+0.040$ & $[-0.021, +0.100]$ & 0.106 & 0.317 \\
MAD & \textbf{Set $\times$ scope interaction} & $\mathbf{+0.060}$ & $\mathbf{[+0.045, +0.075]}$ & \textbf{0.003} & \textbf{0.017} \\
MAD & Weighted vs.\ uniform (tail set)    & $+0.002$ & $[-0.031, +0.035]$ & 0.808 & 0.808 \\
MAD & Weighted vs.\ uniform (weak set)    & $-0.012$ & $[-0.036, +0.013]$ & 0.172 & 0.345 \\
\midrule
MAR20 & Tail arms vs.\ baseline (tail scope)  & $+0.037$ & $[+0.015, +0.058]$ & 0.018 & 0.091 \\
MAR20 & Weak arms vs.\ baseline (weak scope)  & $+0.032$ & $[-0.024, +0.087]$ & 0.133 & 0.336 \\
MAR20 & Set $\times$ scope interaction & $+0.052$ & $[+0.021, +0.083]$ & 0.018 & 0.091 \\
MAR20 & Weighted vs.\ uniform (tail set)    & $+0.014$ & $[-0.008, +0.035]$ & 0.112 & 0.336 \\
MAR20 & Weighted vs.\ uniform (weak set)    & $+0.004$ & $[-0.019, +0.027]$ & 0.519 & 0.519 \\
\bottomrule
\end{tabular}
\end{table}

\section{Discussion}\label{sec:discussion}

\subsection{Engineering guidance: a three-step decision rule}\label{sec:rules}

Read together, the four cells yield an actionable procedure for
practitioners who must decide whether --- and how --- to spend a
synthetic-data budget on a detection system.

\paragraph{Step 1: assess detector headroom before investing in
augmentation}
The largest single improvement observed in this study came not from any
data intervention but from replacing the detector: the pretrained
transformer's \emph{unaugmented} baseline exceeded the best augmented
result of the compact detector by 0.21 \mapm{} on MAD (0.869 vs.\
0.661), and once the stronger detector was in place, every augmentation
we measured was ineffective and resampling was mildly harmful
(Section~\ref{sec:r2}). We stress that this is a cross-cell engineering
observation rather than a randomized comparison --- the two detectors
differ simultaneously in capacity and architecture family (both are
initialized from COCO-pretrained weights;
Section~\ref{sec:detectors}) --- but its practical reading held in every
condition we tested: when deployment constraints permit a stronger
pretrained detector, the model upgrade dominated every augmentation
decision we evaluated. Data-side interventions become worth considering
primarily where a compact detector is mandated, for example by
edge-compute budgets, and the model therefore retains visible headroom.

\paragraph{Step 2: use free resampling as a first attempt --- and as a
diagnostic}
Repeat-factor sampling requires no generation compute and, on the deeply
imbalanced MAD (imbalance ratio 10.5), outperformed every diffusion arm
on every scope (Section~\ref{sec:r3}). Its benefit did not transfer to
the shallowly imbalanced MAR20, where it was null overall and negative
on the weak scope. Because resampling can only re-expose existing
images, it helps precisely when under-exposure of rare classes is the
binding constraint; its failure therefore signals that the constraint
lies elsewhere --- in the diversity of the training distribution rather
than in exposure frequency --- which is the regime where injecting
genuinely new variation through generation remains the available lever.
In practice, RFS functions as a low-cost probe: run it first, keep it if
it helps, and read a null result as evidence that a generative budget
may be worth its price.

\paragraph{Step 3: when generating, select the class set by measured
difficulty; per-class quota tuning is unnecessary}
The interaction result (Section~\ref{sec:r1}) shows that the allocation
signal determines where gains occur, whereas the within-set weighting
axis was null in all four pre-registered comparisons. This simplifies
deployment considerably: the only consequential choice is which classes
to target. Frequency is the conventional default for that choice, but it
is a poor one whenever frequency and difficulty decouple, as they do
here ($r = -0.375$). If the operational goal is to lift worst-case class
performance, the per-class validation AP of the current baseline ---
available at no additional cost in any trained system --- is the direct
signal, and targeting the lowest-AP classes aligns the expected gains
with that goal; within the chosen set, a uniform split of the budget
suffices. The compute cost is moderate: generating and verifying one
thousand 512-pixel background variants took a few GPU-hours on a single
mid-range accelerator in our experiments, of the same order as one
detector training run at this scale.

\subsection{Why background inpainting behaves this way}

Three observations jointly constrain the interpretation. First, the
compact detector's dominant error mode on MAD is missed detection:
baseline recall (0.52--0.59) trails precision (0.62--0.71), and the
largest off-diagonal mass in the confusion structure assigns true
objects to background. Second, gains concentrate on the targeted set yet
consistently spill over, attenuated, to non-targeted classes
(Section~\ref{sec:r1}). Third, all gains vanish under the
higher-capacity pretrained detector (Section~\ref{sec:r2}). These
observations are consistent with a single account: background inpainting
counteracts \emph{background overfitting}. Re-rendering the surroundings
of a geometrically protected object weakens context cues spuriously
correlated with class identity, forcing the detector to rely on object
appearance and thereby recovering missed detections --- most strongly
for classes whose contexts were re-rendered (targeting), positively but
more weakly for classes that share the diversified background statistics
(spillover), and not at all for a detector with the capacity to retain
the broad variation of its pretraining, in which the spurious
object--context coupling is already weak (headroom). The account also
rationalizes the null weighting axis as saturation: once a class
receives enough background variation to break the spurious correlation,
additional variants of the same kind contribute little. Our data contain
direct support for this reading --- the weighted tail arm, whose
smallest realized per-class quota was 16 images against the uniform
arm's 77, achieved statistically indistinguishable tail-scope gains.
Finally, the MAR20 contrast between null resampling and effective
inpainting argues that re-exposure of the protected objects cannot by
itself account for the gains: the synthetic images repeat foregrounds
exactly as resampling does, yet only the arm that also renews the
backgrounds improved performance there. We advance this account as an
interpretation consistent with all of our observations, not as a
demonstrated causal chain.

\subsection{Limitations and threats to validity}

Four limitations bound our claims. (1)~\emph{Replication power.} The
satellite-domain replication reproduces the interaction estimate inside
the original confidence interval, but its six-class scopes leave the
seed-blocked tests underpowered, and the effect does not survive Holm
correction there; we accordingly claim replication of magnitude and
direction, not a second confirmatory significance. (2)~\emph{Set
overlap.} On MAR20 the frequency-defined and difficulty-defined sets
share one of six classes, slightly diluting the dissociation relative to
the fully disjoint MAD design. (3)~\emph{Strong-detector cells: coverage
and attribution.} The transformer cells omit the weak-set arms, so the
headroom finding rests on near-zero deltas of the measured arms rather
than on an interaction test; moreover, the strong-detector comparison
co-varies capacity and architecture family (the pretraining corpus is
shared), so ``headroom'' names their joint effect rather than isolating
a single factor. A full factorial on the strong detector, together with
capacity-matched comparisons within one architecture family, would
resolve both points. (4)~\emph{External validity.} All evidence comes
from one application domain (aircraft), two datasets, one generator
family (Stable Diffusion inpainting, at one budget point per dataset),
and two detector capacity tiers. The condition map should therefore be
read as a demonstrated pattern with identified moderators --- detector
headroom and imbalance depth --- not as a parametric law. Natural next
steps are budget scaling curves, a factorial combination of resampling
with targeted generation to test whether their gains compose, stronger
generators, and replication on broader long-tailed benchmarks.

\section{Conclusion}\label{sec:conclusion}

We asked a question that the growing literature on generative
augmentation leaves implicit: when a fixed budget of synthetic images is
available, which classes should receive it? A $2\times2$ factorial over
the allocation signal and the within-set weighting --- run under a fixed
budget, generator, and verification protocol on two public
military-aircraft benchmarks and two detector families --- gives a
consistent answer. The class set is the decision that matters: gains
land on the targeted set, confirmed by a pre-registered interaction on
the natural-view benchmark and replicated in effect size on the
satellite-view benchmark, while dividing the budget inside the chosen
set uniformly or by priority makes no measurable difference. The choice
of signal is not innocuous, because frequency and measured difficulty
select disjoint classes here; practitioners whose goal is worst-case
class performance should target the measured-weakest classes rather
than the rarest. Both findings, however, are conditional: under a large
pretrained transformer detector every augmentation effect disappears,
and the zero-cost resampling baseline that dominates generation under
deep imbalance turns null --- and locally harmful --- when the imbalance
is shallow. The resulting decision rule --- assess headroom, probe with
free resampling, then generate for measured weakness with a uniform
split --- is immediately actionable, and the evaluation protocol that
produced it (budget control, pre-registered contrasts, seed-blocked
tests) is reusable wherever augmentation claims are made. Extending the
factorial to strong detectors, tracing budget scaling curves, and
testing whether resampling and targeted generation compose are the
natural next steps.

\bibliographystyle{elsarticle-harv}
\bibliography{refs}

\end{document}
